\begin{theorem}[Linear expected contact time across every cut]
\label{thm:c2readycontact}
Consider the adopted algorithm with $c=\ell=2$ under its
independent uniform global initiator scheduler. Fix a
nontrivial partition $V=S\sqcup G$, and suppose some
cache in $S$ contains an identity $r\in G$.
Let $\tau$ be the first exchange crossing this cut,
in either direction. Then, uniformly over arbitrary
finite initial ages and all cache contents,
\[
 \mathbb E_x\tau\le7n,\qquad
 \Pr_x\{\tau>14nk\}\le2^{-k}\quad(k\ge0).
\]
No controller, frozen scheduling rule, or prescribed
sequence of initiators is used.
\end{theorem}

\begin{proof}
Until the first crossing, all events are internal to
$S$ or internal to $G$. Let $N$ count the caches in
$S$ containing identity $r$. It cannot decrease before
$\tau$. Indeed, in any $S$-internal exchange the owner
$r$ is not an endpoint. If both endpoints hold $r$,
the request's extra and the complete reply protect it
at both endpoints. If just one holds $r$, it is
forwarded to and protected at the opposite endpoint,
whether or not the original location survives.
Events internal to $G$ do not change the $S$ caches.
Thus $N\ge1$ throughout the pre-crossing trajectory.

Call a copy \emph{ready} when its cache is full and
has the form
\[
 C_q=\{r^A,s^B\},\qquad A\ge B.
\]
Let $R$ count such caches in $S$. Their initiators
select $r$ with conditional probability at least
$1/2$. Consequently, if $R>0$, the conditional
probability $a$ of a crossing at the next global
event satisfies
\[
 a\ge\frac{R}{2n}.
\]

We need two elementary facts about losses of readiness:
a noncrossing event decreases $R$ by at most one,
and its conditional probability $d$ of doing so obeys
\[
 d\le\frac{2R}{n}.
\]
Here is the complete charging argument. Only the
initiator $p$ and responder $q$ can change.
If $p$ is ready and does not cross, its other head
$q$ must be tied with $r$. Its request contains
fresh $p^0$ and $r$, so the responder ends with
a full ready cache. At worst the initiator loses
its ready status; even if the responder was already
ready, the total loss is at most one.
Charge all these events to the clocks of the
$R$ ready initiators, of total probability $R/n$.

It remains to inspect a ready responder
$C_q=\{r^A,s^B\}$ when the initiator was not ready.
If the request includes $r$, the responder ends
with $\{p^0,r^C\}$ and stays ready.
If it does not include $r$ and $p\ne s$, the
complete reply contains two usable identifiers
$r,s$. A singleton initiator receives both and
is ready. A full initiator either discards its
novel extra to receive both, or its extra is
$s$, whose minimum-merged age is at most
$B\le A$. Again readiness transfers to the
initiator whenever it is lost at the responder.
The sole remaining possibility is $p=s$:
the reply entry $s$ is self-filtered, and the
initiator's retained extra may be older than $r$.
Thus a loss not already charged to a ready
initiator requires activation of the ready
responder's unique other cached identity $s$.
Union-bounding these at most $R$ clocks costs
at most another $R/n$. The same cases prove
that the total loss is never larger than one.

Use the nonnegative potential
\[
 \Psi(R)=n\bigl(2+5(4/5)^R\bigr)
\]
before a crossing, and set its value to zero after
one. Put $q_*=4/5$. If $R>0$, readiness gains
can only decrease this potential, a loss of one
costs exactly $nq_*^{R-1}$, and a crossing removes
$n(2+5q_*^R)$. Its conditional drift is therefore
at most
\[
 \begin{split}
 \mathbb E[\Delta\Psi\mid\text{present state}]
 &\le-a n(2+5q_*^R)+d n q_*^{R-1}\\
 &\le-\frac R2(2+5q_*^R)+2R q_*^{R-1}
 =-R\le-1.
 \end{split}
\]

If $R=0$, some $r$ holder still exists because
$N\ge1$. Select any such holder mathematically
for a hazard bound; this does not alter the
scheduler. With probability at least $1/n$ it
is the next initiator. A singleton holder
contacts $r$ and crosses. A full nonready
holder has its other entry strictly older
than $r$ and selects that entry. If this
crosses, we are done. Otherwise its request
forces the responder to cache
$\{p^0,r^C\}$ with $C\ge0$, creating a ready
copy. Thus with probability at least $1/n$
the potential falls from $7n$ to at most
$6n$ (or zero); other events cannot increase
it above its value at $R=0$. Consequently
its conditional drift is again at most $-1$.

For any integer $t\ge0$, summing this drift
until $\tau\wedge t$ and using nonnegativity
gives
\[
 \mathbb E_x(\tau\wedge t)
 \le\Psi(R_0)\le7n.
\]
Monotone convergence proves
$\mathbb E_x\tau\le7n$, with no assumption
of finite expected contact time in advance.

Markov's inequality gives
$\Pr_x\{\tau>14n\}\le1/2$ uniformly.
Before a crossing, $N\ge1$ still holds,
so the same estimate applies at every
$14n$-event block boundary.
The Markov property proves the stated tail.
\end{proof}

\begin{corollary}[Ordinary temporal connectivity in quadratic time]
\label{cor:c2readytemporal}
Start from any weakly connected overlay. Let
$\mathcal I_t$ be the undirected graph of pairs
that have exchanged during the first $t$ events,
and let $\tau_{\rm conn}$ be its connectivity time.
Then
\[
 \mathbb E_x\tau_{\rm conn}\le7n(n-1),
 \qquad
 \Pr_x\{\mathcal I_t\text{ is disconnected}\}
 \le(2^{n-1}-1)\,
       2^{-\lfloor t/(14n)\rfloor}.
\]
Thus
\[
 \Pr_x\!\left\{
 \tau_{\rm conn}>
 14n\left\lceil n-1+\log_2(1/\varepsilon)\right\rceil
 \right\}\le\varepsilon.
\]
These are communication bounds; they do not
assert uniform peer sampling at that time.
\end{corollary}

\begin{proof}
For each fixed unordered nontrivial cut, initial
weak connectivity supplies a descriptor across it.
Orient the cut so its tail is in $S$ and its
head $r$ is in $G$. The theorem bounds its
first observed crossing. There are
$2^{n-1}-1$ unordered nontrivial cuts; a union
bound proves the tail.

For the expectation, whenever the observed
interaction graph is disconnected, select one
of its components and keep its vertex set
fixed until the next crossing of that cut.
Conditional on the present overlay, the theorem
bounds the expected wait by $7n$.
This crossing merges observed components,
and at most $n-1$ such phases are needed.
Summing their conditional mean bounds gives
the claimed expectation.
\end{proof}

